\section{Lagrange Multiplier dalam Markowitz-QUBO}\label{appendix-c}

\hypertarget{definisi-mengapa-a-disebut-bobot-penalti-lagrange-multiplier}{%
\subsection{\texorpdfstring{Definisi: Mengapa \(A\) disebut Bobot Penalti (Lagrange Multiplier)?}{Definisi: Mengapa A disebut Bobot Penalti (Lagrange Multiplier)?}}\label{definisi-mengapa-a-disebut-bobot-penalti-lagrange-multiplier}}

Dalam konteks \textbf{QUBO (Quadratic Unconstrained Binary Optimization)}, kata kuncinya adalah ``\emph{Unconstrained}'' (Tanpa Kendala). Algoritma kuantum seperti QAOA atau mesin \emph{annealing} hanya bisa meminimalkan fungsi energi tanpa batasan eksternal.

Namun, dalam dunia nyata, kita memiliki batasan (seperti: ``harus memilih tepat \(K\) aset''). Untuk memasukkan batasan ini ke dalam fungsi yang ``tanpa batasan'', kita menggunakan \textbf{Metode Penalti}.

\begin{itemize}
\tightlist
\item
  \textbf{\(A\) sebagai Bobot Penalti}: Menentukan seberapa besar ``hukuman'' (peningkatan energi) yang diberikan jika solusi melanggar batasan.
\item
  \textbf{\(A\) sebagai Lagrange Multiplier}: Secara fungsional bertindak seperti pengali \(\lambda\) dalam kalkulus, yang menghubungkan fungsi objektif utama dengan fungsi kendala.
\end{itemize}

\hypertarget{hubungan-dengan-metode-lagrange-standar}{%
\subsection{Hubungan dengan Metode Lagrange Standar}\label{hubungan-dengan-metode-lagrange-standar}}

Suku penalti tersebut adalah adaptasi dari \textbf{Metode Lagrange Multiplier} untuk optimasi diskrit. Dalam optimasi kontinu standar:

\[\min f(x) \quad \text{dengan kendala} \quad g(x) = 0\]

Dirumuskan menjadi fungsi Lagrangian: \[\mathcal{L}(x, \lambda) = f(x) + \lambda g(x)\]

Dalam QUBO, kita menggunakan variasi yang disebut \textbf{Penalty Method} (bagian dari \emph{Augmented Lagrangian Method}). Perbedaannya, kita menguadratkan kendalanya: \(A(g(x))^2\).

\textbf{Mengapa dikuadratkan?} 1. \textbf{Selalu Positif}: Jika \(g(x) \neq 0\) (baik positif atau negatif), nilai \((g(x))^2\) akan selalu positif, sehingga menambah energi (memberi penalti). 2. \textbf{Kesesuaian QUBO}: Bentuk kuadrat memastikan hubungan antar variabel tetap dalam orde dua (kuadratik), yang merupakan syarat format QUBO.

\hypertarget{penurunan-rumus-penalti}{%
\subsection{Penurunan Rumus Penalti}\label{penurunan-rumus-penalti}}

Berikut adalah langkah transformasi dari kendala menjadi suku penalti:

\begin{enumerate}
\def\labelenumi{\arabic{enumi}.}
\tightlist
\item
  \textbf{Kendala Asli}: \(\sum_{i=1}^N x_i = K\) (Jumlah aset yang dipilih harus sama dengan \(K\)).
\item
  \textbf{Bentuk Normal Kendala}: \(g(x) = \sum_{i=1}^N x_i - K = 0\).
\item
  \textbf{Konstruksi Penalti Kuadratik}: Agar nilai minimum fungsi tepat berada di titik di mana \(\sum x_i - K = 0\), kita kuadratkan selisihnya: \[ P(x) = A \cdot (g(x))^2 = A \left( \sum_{i=1}^N x_i - K \right)^2 \]
\end{enumerate}

Jika \(\sum x_i = K\), maka \((K - K)^2 = 0\) (Tanpa penalti). Jika terjadi pelanggaran, energi naik sebesar \(A \times (\text{error})^2\).

\hypertarget{fungsi-umum-lagrange-multiplier}{%
\subsection{Fungsi Umum Lagrange Multiplier}\label{fungsi-umum-lagrange-multiplier}}

Fungsi utama Lagrange Multiplier (\(A\)) adalah sebagai \textbf{penyeimbang antara dua tujuan yang saling bertentangan}:

\begin{enumerate}
\def\labelenumi{\arabic{enumi}.}
\tightlist
\item
  \textbf{Optimalitas (Fungsi Objektif)}: Mencari nilai sekecil atau sebesar mungkin (misal: risiko terkecil atau keuntungan terbesar).
\item
  \textbf{Validitas (Kendala)}: Memastikan solusi tersebut memenuhi syarat yang ditetapkan (misal: jumlah aset pas).
\end{enumerate}

\textbf{Analogi Sederhana:} Bayangkan Anda sedang berbelanja. * \textbf{Fungsi Objektif}: Membeli barang sebanyak mungkin (memuaskan keinginan). * \textbf{Kendala}: Uang di dompet terbatas. * \textbf{Lagrange Multiplier (\(A\))}: Adalah ``konsekuensi'' jika Anda \emph{over-budget}. Jika \(A\) terlalu kecil, Anda mungkin nekat berutang (melanggar kendala). Jika \(A\) sangat besar, Anda akan sangat hati-hati agar pengeluaran pas dengan dompet, meskipun keinginan belum sepenuhnya terpenuhi.

Dalam QUBO, nilai \(A\) harus dipilih agar cukup besar sehingga kendala tidak dilanggar, tetapi tidak terlalu besar sehingga menutupi detail dari fungsi risiko dan return yang ingin kita optimalkan.

\hypertarget{strategi-menentukan-nilai-a}{%
\subsection{\texorpdfstring{Strategi Menentukan Nilai \(A\)}{Strategi Menentukan Nilai A}}\label{strategi-menentukan-nilai-a}}

Menentukan nilai \(A\) yang tepat adalah salah satu tantangan utama dalam QUBO. Jika \(A\) terlalu kecil, algoritma akan menghasilkan solusi yang melanggar batasan. Jika \(A\) terlalu besar, ``lanskap energi'' menjadi terlalu curam sehingga sulit bagi algoritma untuk menemukan solusi optimal global.

Berikut adalah beberapa pendekatan untuk mendapatkan nilai \(A\):

\hypertarget{a.-pendekatan-teoretis-upper-bound}{%
\subsubsection{Pendekatan Teoretis (Upper Bound)}\label{a.-pendekatan-teoretis-upper-bound}}

Secara matematis, \(A\) harus lebih besar daripada perubahan maksimum yang mungkin terjadi pada fungsi objektif utama (risiko dan return) jika satu variabel berubah. \[ A > \max | \Delta (\text{Risiko} - \text{Return}) | \] Tujuannya adalah memastikan bahwa keuntungan dari ``curang'' (melanggar batasan) tidak pernah lebih besar daripada kerugian akibat hukuman penalti.

\hypertarget{b.-rule-of-thumb-praktis}{%
\subsubsection{Rule of Thumb (Praktis)}\label{b.-rule-of-thumb-praktis}}

Seringkali, praktisi menggunakan nilai \(A\) yang setara atau sedikit lebih besar dari nilai koefisien terbesar dalam matriks kovarians atau return. - Jika \(\Sigma_{ij}\) dan \(\mu_i\) berada di rentang \([0, 1]\), nilai \(A\) sering kali dimulai dari \(1\) hingga \(10\).

\hypertarget{c.-pencarian-grid-grid-search-hyperparameter-tuning}{%
\subsubsection{Pencarian Grid (Grid Search / Hyperparameter Tuning)}\label{c.-pencarian-grid-grid-search-hyperparameter-tuning}}

Dalam pendekatan ini, \(A\) diperlakukan seperti hyperparameter dalam machine learning: 1. Pilih serangkaian nilai (misal: \(0.1, 1, 10, 100\)). 2. Jalankan optimasi untuk setiap nilai. 3. Pilih nilai \(A\) terkecil yang masih memberikan solusi valid (memenuhi \(\sum x_i = K\)).

\hypertarget{d.-penalti-dinamis-annealing-schedule}{%
\subsubsection{Penalti Dinamis (Annealing Schedule)}\label{d.-penalti-dinamis-annealing-schedule}}

Beberapa algoritma memulai dengan nilai \(A\) yang kecil agar sistem dapat mengeksplorasi ruang solusi secara bebas pada awal proses, kemudian secara bertahap meningkatkan \(A\) untuk ``memaksa'' sistem masuk ke wilayah solusi yang valid di akhir proses.

\hypertarget{e.-normalisasi-fungsi-objektif}{%
\subsubsection{Normalisasi Fungsi Objektif}\label{e.-normalisasi-fungsi-objektif}}

Agar pemilihan \(A\) lebih konsisten, seringkali suku risiko dan return dinormalisasi terlebih dahulu ke rentang \([0, 1]\). Dengan demikian, nilai \(A\) dapat ditentukan secara lebih intuitif (misal: \(A=1.5\) pasti lebih dominan daripada fungsi objektif yang sudah dinormalisasi).
